% 原文 pp.261–268（印刷 pp.1379–1386）

\section{Lecture 19：四维 $\mathcal{N}=2$ 超对称 Yang--Mills 理论（三）：拓扑应用}

\begin{keytakeaways}
本讲把前两讲在 Minkowski 四维空间上得到的 $\mathcal{N}=2$ 超对称 Yang--Mills 低能动力学，搬到一般的紧致定向 Riemann 四维流形 $X$ 上，为 Donaldson 不变量寻找 Seiberg--Witten 式的表述。
定性上的分水岭是单圈 $\beta$ 函数 $\beta_1(g)=-\frac{g^3}{8\pi^2}\bigl(4h-c_2(\rho)\bigr)$ 的符号：为负则理论渐近自由、是好的基本理论，为零则 $\beta\equiv0$、$\tau$ 是常数且理论具有 $\Gamma\subset SL(2,\mathbb{Z})$ 的对偶对称性，为正则理论在紫外不自由、只能充当红外极限的候选者。
数学工具是扭变（twisting）：把 $Spin(4)=SU(2)_+\times SU(2)_-$ 中的 $SU(2)_+$ 与 $\mathcal{N}=2$ 的 $R$ 对称性 $SU(2)_R$ 取对角，得到 $SU(2)'$，于是超对称生成元 $Q^i_\alpha$ 与 $\bar Q_{j\dot\alpha}$ 都落入 $Spin(4)$ 的矢量表示，从而可在任意定向 Riemann 四维流形 $X$ 上整体定义。
扭变的几何后果是：在分裂超流形 $\Lambda^0TX\oplus TX\oplus\Lambda^2_+(TX)$ 上存在一个处处非零的整体实向量场 $Q=\partial_{\mathbf{R}^{\rm odd}}+\theta^i\partial_{i,{\rm even}}$，满足 $Q^2=0$ 与 $\{Q,K\}=d$；而由 $\{Q,\mathcal{S}\}=T$ 可得 $\{Q,T\}=0$，故所关心的关联函数与度规无关。
主要结论是：由 Chern 形式与手征标量 $\phi$ 造出的 $Q$-闭算符 $\mathcal{O}_P^{(n)}=K^nP(\phi)$ 正是 Donaldson 类在此框架中的对应物，其期望值等于各拓扑类型的模空间 $\mathcal{M}_E$ 上 $\mu_P(\Sigma)$ 的杯积积分；而扭变之后矢量多重态的全部场都化为微分形式，因而在任意闭定向四维流形上都有定义。
\end{keytakeaways}

% ===== 原文 p.261 =====
% 原文标题页：LECTURE 19 / $N=2$ Super-symmetric Yang-Mills Theories in Dimension Four: Part 3, Topological Applications

\subsection{四维 $\mathcal{N}=2$ 超对称规范理论概览}
% 原文 §19.1：A survey of $N=2$ supersymmetric gauge theories in dimension four

让我们先对 $\mathcal{N}=2$ 超对称规范理论的各类情形作一简要概述\footnote{原文页脚：\emph{Notes by John Morgan}——本讲笔记由 John Morgan 整理。}。我们必须指定一个规范群 $G$，它将是一个紧致连通群。为得到 $\mathcal{N}=2$ 超对称，我们的场中包含一个矢量多重态（vector multiplet），它由规范场及其超伙伴场构成，取值于主 $G$-丛的伴随丛（adjoint bundle）。我们也可以有超多重态（hypermultiplet）。设 $n$ 为超多重态的数目。$n$ 个自由超多重态的全局对称群是辛群 $Sp(n)$（即 $n\times n$ 四元数矩阵构成的酉群）。以 $\mathcal{N}=2$ 不变的方式对理论作规范化（gauging），可以通过规范化 $Sp(n)$ 的任意所需子群来实现。因此一般而言，超多重态位于由表示 $\rho:G\to Sp(n)$ 与主 $G$-丛相伴的向量丛中。

对理论的定性性质而言最重要的不变量是 $\beta$ 函数，尤其是单圈 $\beta$ 函数。单圈 $\beta$ 函数的公式为
\[
\beta_1(g)=-\frac{g^3}{8\pi^2}\bigl(4h-c_2(\rho)\bigr),
\]
其中 $h$ 是群 $G$ 的对偶 Coxeter 数（dual Coxeter number），$c_2(\rho)$ 是表示 $\rho$ 的二次 Casimir 的迹，其归一化使得 $SU(2)$ 的定义表示具有 $c_2$ 等于一。当然，$2h$ 正是 $c_2(\text{adjoint})$，即 $G$ 的伴随表示在其李代数上的二次 Casimir。有三种可能性：
\begin{itemize}
\item 单圈 $\beta$ 函数为零。
\item 单圈 $\beta$ 函数为负。
\item 单圈 $\beta$ 函数为正。
\end{itemize}

% [Tutor Note: 此处 $c_2$ 的归一化是核对本讲一切公式的第一道关口。作者取「$SU(2)$ 定义表示的 $c_2$ 等于一」，即 $c_2(\rho)=2T(\rho)$（国内教材常用的指标 $T$ 以 $T(\text{fund})=1/2$ 归一化）；只有在这个约定下才有「$2h=c_2(\text{adjoint})$」，这也正是 $h$ 被称为\emph{对偶} Coxeter 数的原因。最容易出错的是括号里的 $4h$ 而非 $2h$：规范多重态贡献 $2c_2(\text{adj})$，而 $c_2(\text{adj})=2h$，故为 $4h$。可自行校验：下一页（p.262）说 $c_2(\rho)=4h$ 对应「伴随表示复化的两份」，而每份伴随的 $c_2$ 是 $2h$，两份恰为 $4h$——若把 $4h$ 误抄成 $2h$，这个 $\mathcal{N}=4$ 判据立刻失效，且纯 $SU(2)$ 的 $\beta_1$ 会差一个因子 $2$。]

由于
\[
\tau=\frac{\theta}{2\pi}+\frac{4\pi i}{g^2},
\]

% ===== 原文 p.262 =====
我们有展开
\[
\mu\frac{d}{d\mu}(\tau)=\text{constant}+\sum_{\ell\ge1}C\cdot g^{2\ell}+\sum c_{n,\ell}\,g^{2\ell}\exp(2\pi i n\tau),
\]
其中常数项由单圈 $\beta$ 函数算出，$g$ 的更高次幂来自更高圈的贡献，而指数项来自瞬子修正（instanton correction，它们在微扰论中当然是看不到的）。由超对称性，$\mu\frac{d}{d\mu}(\tau)$ 是全纯函数，因此该展开中 $g$ 的正幂次项恒为零。于是，若单圈 $\beta$ 函数为零，则 $\beta$ 在微扰论中为零。在这种情况下，人们相信 $\beta\equiv0$。这一点尚未被严格证明，但显然，当单圈 $\beta$ 函数为零时，我们在前两讲中分析过的那类低能有效场论，只有在 $\beta$ 函数恒等于零时才有意义。（相反，单圈 $\beta$ 函数非零的 $\mathcal{N}=2$ 理论确实有非微扰的瞬子修正，正如我们在纯 $SU(2)$ 理论中所见。）

% [Tutor Note: 「由超对称性 $\mu\frac{d}{d\mu}\tau$ 是全纯函数，故 $g$ 的正幂次项恒为零」这句跳过了关键一步：全纯性是就 $\tau$ 而言的（在 $\mathcal{N}=2$ 中 $\tau$ 是手征超场的 $F$ 项系数，重整化只允许它作全纯移动），而 $\tau$ 与 $g$ 的关系是 $g^2=4\pi i/(\tau-\theta/2\pi)$；把 $\mu\frac{d\tau}{d\mu}$ 写成 $g$ 的幂级数时，$\ell\ge1$ 的每一项都隐含 $\ln\Lambda$ 型的非全纯依赖，故必为零。换算关系本身值得记住并用来验算：由 $\tau$ 的定义直接求导得 $\mu\frac{d\tau}{d\mu}=-8\pi i\,g^{-3}\beta(g)$，代入 $\beta_1=-\frac{g^3}{8\pi^2}(4h-c_2(\rho))$ 得常数项 $i(4h-c_2(\rho))/\pi$；反之若把 $\tau$ 写成常见的 $\theta/2\pi+i/g^2$（少一个 $4$），这里就会整体差因子 $4$。另注：指数项中 $n$ 的求和未写范围，实为 $n\ge1$，且瞬子展开的收敛性要求 $\operatorname{Im}\tau>0$（即 $g^2>0$），这正是 $\tau$ 取上半平面的原因。]

$\mathcal{N}=2$ 理论中 $c_2(\rho)=4h$ 的一个重要特例，是超多重态表示 $\rho$ 为伴随表示复化的两份这一情形。\footnote{任一秩为 $n$ 的实表示 $W$——这里 $W$ 是伴随表示，$n$ 是群的维数——都可以看作 $SO(n)$ 中的一个表示。$W\otimes_{\mathbb{R}}\mathbb{C}$ 的两份可以看作辛群 $Sp(n)$ 中的一个表示。为简洁起见，我们将把「$W$ 的复化的两份」称作「$W$ 与四元数作（$\mathbb{R}$ 上的）张量积」，其中四元数就简单地理解为带非退化斜形式的 $\mathbb{C}^2$。} 这个理论实际上具有 $\mathcal{N}=4$ 超对称。利用这一增强的超对称性可以证明，$\beta$ 函数恒为零。

当 $\beta\equiv0$ 时，量子理论与经典理论一样具有一个定义良好的常数 $\tau$。人们相信，在这种情况下存在 $SL(2,\mathbb{Z})$ 中一个有限指数的子群 $\Gamma$，使得理论在 $\Gamma$ 的忠实作用下不变。当然，当 $\beta\ne0$ 时理论有质量标度，$\tau$ 不再是常数。对于这些理论，人们会预期（正如我们在前两讲中对纯 $SU(2)$ 理论所见）：理论并不存在精确的对偶对称性，反倒是各种可能的低能描述之间的对偶，对于理解低能物理十分重要。

当单圈 $\beta$ 函数为负时，$\beta$ 函数为负，至少对足够高的能标如此。（我们上次研究的纯 $SU(2)$ 理论即属此类。）在这种情况下，我们有一个渐近自由（asymptotic freedom）的理论，它是一个好的基本理论。也就是说，我们有一个提法良好的\emph{问题}（即描述该理论的低能有效版本）。我们在上一讲中对 $G=SU(2)$ 或 $SO(3)$、$\rho$ 为平凡表示的情形解决了这个问题。

当单圈 $\beta$ 函数为正时，我们没有提法良好的问题，因为理论不是渐近自由的。毋宁说，这些理论是那些单圈 $\beta$ 函数为负的理论所提出的问题的候选解。更确切地说，一个 $\beta$ 函数为正的理论 $T_{\rm mac}$ 之所以是候选解，是因为它的低能有效理论提法良好且在红外自由。它是 $\beta$ 为负的理论 $T_{\rm mic}$ 所提问题的解，如果 $T_{\rm mic}$ 在红外流向 $T_{\rm mac}$，即在长距离（或等价地，低能量）下两个理论的关联函数相互收敛。把问题与答案配对、弄清楚哪些 $T_{\rm mac}$ 可以作为给定 $T_{\rm mic}$ 的红外极限出现，这个一般问题尚未解决，尽管在许多有趣的情形中已有大量结果。

% ===== 原文 p.263 =====
我们今天讨论的规范理论，是由量子化经典物理而得到的 Lagrange 模型。至少顺带一提：四维中还存在更为奇特的 $\mathcal{N}=2$ 理论，其中一些来自弦论。并非所有这些理论都有已知的 Lagrange 表述。无论它们是否有 Lagrange 量，这些 $\mathcal{N}=2$ 超对称理论在扭变之后也应当给出四维流形的不变量。那些没有 Lagrange 量的理论，可能给出这样的四维流形不变量：它们没有以某个微分方程之解的模空间为基础的经典描述。

\subsection{从 Minkowski 空间到紧致 Riemann 四维流形}
% 原文 §19.2：From Minkowski space to a compact Riemannian four-manifold

让我们固定一个 $\mathcal{N}=2$ 超对称渐近自由规范理论（即 $\beta<0$ 的那一个）。例如，若我们想研究 Donaldson 理论，就取纯 $\mathcal{N}=2$ 超对称 $SU(2)$ 规范理论。但在一段时间内，我们可以更一般些。到目前为止我们所讨论的理论，都位于 Minkowski 四维空间上。也就是说，我们写下的作用量涉及在 Minkowski 空间上对一个由该空间上各种场构成的 Lagrange 函数作积分。我们希望从它过渡到一个定义在紧致 Riemann 四维流形上的理论。这一过渡分两步完成——i) 从 Minkowski 四维空间转动到欧氏四维空间；ii) 整体化（globalize）到紧致 Riemann 四维流形。

让我们回忆去年秋天学到的、如何把理论 Wick 转动到欧氏四维空间。我们的理论由一族带有算符乘积展开（operator product expansion，O.P.E.）的局部算符构成。理论的一个实现（realization）给出这些算符在 Minkowski 四维空间各点处关联函数的值。但在 Minkowski 四维空间复化之后的一个开子集上，O.P.E. 与关联函数可由解析延拓定义。这个开子集包含欧氏四维空间。我们就是这样通过 Wick 转动把理论定义在欧氏四维空间上（例如见 [I-Axiomatics] 的第 2 讲）。在我们的情形中，所处理的是（超对称的）作用量
\[
S=\int_{M^4}\mathcal{L},
\]
它们是 Lagrange 量 $\mathcal{L}$ 的积分，而 $\mathcal{L}$ 是涉及底流形 Minkowski 空间上诸场 $\phi$ 的解析表达式。这些表达式延拓为复化 Minkowski 空间上的全纯表达式，然后可以限制到欧氏子空间上，产生欧氏 Lagrange 量 $\mathcal{L}_E$。欧氏作用量定义为
\[
S_E=\int_{\mathbb{R}^4}\mathcal{L}_E.
\]
Minkowski 空间中关联函数的路径积分为
\[
\int\mathcal{D}\phi\,\mathcal{O}_1(x_1)\cdots\mathcal{O}_t(x_t)\,e^{iS},
\]
而在这一过程中它 Wick 转动为路径积分
\[
\int\mathcal{D}\phi\,\mathcal{O}_1(x_1)\cdots\mathcal{O}_t(x_t)\,e^{-S_E},
\]

% ===== 原文 p.264 =====
计算的是欧氏空间中的关联函数。即使人们关心的是 Minkowski 空间，计算也常常在欧氏空间中进行，因为那里有一些技术上的简化。关于二者关系的更多细节，见 [I-Classical Fields] 第 7 章。

\subsubsection{通过扭变实现整体化——超对称代数情形}
% 原文 §19.2.1：Globalizing by twisting – the case of the supersymmetry algebra

现在让我们考察把该量子场论放到定向 Riemann 四维流形 $X$ 上的问题。我们的超对称理论是在四维空间上由一个作用量 $S$ 定义的，$S$ 是局部场的 Lagrange 函数的积分；这些场在 $Spin(4)$ 下按标量、旋量与矢量变换。当然，在 Minkowski 或欧氏四维空间上，有与这些表示相伴的丛。为了在四维流形 $X$ 上写下「同一个」作用量，我们需要能够在该流形上写下超对称 Lagrange 量。这意味着，Lagrange 量中出现的所有场都必须在 $X$ 上整体定义（作为丛的截面）。就目前所呈现的理论而言，一般情形下并非如此，因为相关的旋量丛可能不在 $X$ 上整体存在。事实上，Riemann 四维流形上的「物理」理论（即真正把理论与广义相对论耦合时会用到的那种）要求 $X$ 是 spin 流形，并且一般会破坏掉全部超对称。我们此前已经看到（第 14 讲 §3），对拓扑应用如何处理这个问题——我们用一个适当的同态把自旋群 $Spin(4)$ 映入 $R$ 对称群，以此扭变（twist）理论。我们将会看到，适当的扭变既使我们能在任意定向 Riemann 四维流形上写下 Lagrange 量，又使我们能保住超对称的一个关键部分。

把 $\Lambda^2T^*\mathbb{R}^4$ 分解为自对偶与反自对偶两部分，这由群 $Spin(4)$ 分解为 $SU(2)_+\times SU(2)_-$ 所覆盖。幸运的是，对 $\mathcal{N}=2$，$R$ 对称群是 $U(2)_R=SU(2)_R\times U(1)_R$。（$U(1)_R$ 常常并不存在，例如当超多重态有非零裸质量时；当 $\beta\ne0$ 时它也有反常。但即便如此，我们仍以 $SU(2)_R$ 作为 $R$ 对称群。）有了这个 $R$ 对称群，扭变才有可能，因为存在从 $Spin(4)$ 到 $SU(2)_R$ 的非平凡同态。（对 $\mathcal{N}=1$，$R$ 对称性是 $U(1)_R$，而在四维中不存在扭变的可能性。）回忆 [I-Supersolutions] 第 10 章：超对称生成元是 $\mathbb{R}^{4|8}$ 上的右不变向量场，记为 $Q^i_\alpha$ 与 $\bar Q_{j\dot\alpha}$，其中 $1\le i,j,\alpha\le2$。这两组生成元都在 $i,j$ 指标上按标准方式受 $SU(2)_R$ 作用。（$\alpha$ 与 $\dot\alpha$ 是 $SU(2)_+\times SU(2)_-$ 的旋量指标。）自旋群与 $SU(2)_R$ 的作用归纳在下表中。

\begin{center}
\begin{tabular}{|c||c|c|c|}
\hline
\multicolumn{1}{|c||}{} & \multicolumn{3}{c|}{各群下的变换}\\
\hline
\multicolumn{1}{|c||}{荷} & $SU(2)_+$ & $SU(2)_-$ & $SU(2)_R$\\
\hline
$Q^i_\alpha$ & $\frac12$ & $0$ & $\frac12$\\
\hline
$\bar Q_{j\dot\alpha}$ & $0$ & $\frac12$ & $\frac12$\\
\hline
\end{tabular}
\end{center}

% [Tutor Note: 这张表的指标位置是全书最容易抄错的地方之一。$Q^i_\alpha$ 的\emph{上}指标 $i$ 是 $SU(2)_R$ 的指标（$i=1,2$），\emph{下}指标 $\alpha$ 是 $SU(2)_+$ 的旋量指标；$\bar Q_{j\dot\alpha}$ 则相反：$j$ 在下、且带点的 $\dot\alpha$ 属于 $SU(2)_-$。扭变时被「吃掉」的正是 $i$ 与 $\alpha$ 这一对（对角化 $SU(2)_+$ 与 $SU(2)_R$），所以若把两者写反，下一页 $Q^i_\alpha$ 分解为 $(0,0)\oplus(1,0)$ 的结论会变成另一个表示，整个扭变就错了。另外，$SU(2)_-$ 的那一列对两个生成元分别给出 $0$ 与 $\frac12$，这恰是 $\bar Q$ 带点指标 $\dot\alpha$ 的体现——表中 $SU(2)_-$ 列非零的行正是 $\bar Q$ 所在的行。]

我们定义 $SU(2)'$ 为 $SU(2)_+\times SU(2)_R$ 在对这两个群的显然认同（identification）之下的对角子群。仍然在 $\mathbb{R}^4$ 上工作，我们将关心

% ===== 原文 p.265 =====
Poincaré 群的一个作用，其中通常的转动群 $Spin(4)$ 被 $SU(2)'\times SU(2)_-$ 取代。由表中的信息可以推出，在 $SU(2)'\times SU(2)_-$ 下，$Q^i_\alpha$ 按表示
\[
\Bigl(\frac12\otimes\frac12,0\Bigr)=(0,0)\oplus(1,0)
\]
变换。这意味着，$Q^i_\alpha$ 的扭变形式分解为一个复值向量场，与一个取值于复自对偶二形式的向量场。类似地，$\bar Q_{j\dot\alpha}$ 按同一群的表示
\[
\Bigl(\frac12,\frac12\Bigr)
\]
变换，从而是切向量，于是 $\bar Q_{j\dot\alpha}$ 成为取值于复化切丛的向量场。因此，这两个表示都是 $Spin(4)$ 的矢量表示，从而在任何定向 Riemann 四维流形上都整体存在。另外请注意，这些表示自然地是实表示（尽管原来的表示并不是）。这意味着在 Riemann 流形上我们可以取实场。

% [Tutor Note: 原文此处作「\emph{If} follows from the information in the table that…」，显系「\emph{It} follows…」之误（排印脱字），译文按正确读法处理。数学上这一步的实质是 $SU(2)$ 的张量积分解 $\frac12\otimes\frac12=0\oplus1$：$Q^i_\alpha$ 原本在 $SU(2)_+\times SU(2)_R$ 下按 $(\frac12,\frac12)$ 变换，对角化之后该表示限制为 $0\oplus1$，其中单态 $(0,0)$ 就是后来那个整体存在的标量型生成元 $Q$，三重态 $(1,0)$ 则给出取值于自对偶二形式的向量场；而 $\bar Q$ 的 $(\frac12,\frac12)$ 在对角化后仍是一个「矢量」，因为 $\frac12\otimes\frac12$ 中恰含一个自旋 $1$ 分量、且其 $SU(2)_-$ 指标 $\dot\alpha$ 保持不动（故为反自对偶侧）。]

用超空间的语言说，扭变并过渡到一般四维流形，意味着把超流形 $\mathbb{R}^{4|8}$（即 $\mathbb{R}^4$ 的超对称扩充）替换为一个分裂超流形（split supermanifold），其偶部是四维流形 $X$，奇部是宇称反转的向量丛 $\Lambda^0TX\oplus TX\oplus\Lambda^2_+(TX)$。在这个超流形上，我们有如下定义的一个整体实向量场 $Q$。设我们在 $X$ 的某个开子集上有局部坐标 $(x^1,\dots,x^4)$。令 $\theta^{\mathbf{R}}$ 为奇的 $\Lambda^0TX$ 方向上的自然坐标，$\theta^i$ 为由 $x^i$ 决定的奇的 $TX$ 方向上的坐标。我们以 $\partial_{\mathbf{R}^{\rm odd}}$ 与 $\partial_{i,{\rm odd}}$ 记这些奇方向上相应的向量场。最后，$\partial_{i,{\rm even}}$ 是沿 $x^i$ 方向的通常偏导数。我们写
\[
Q=\partial_{\mathbf{R}^{\rm odd}}+\theta^i\partial_{i,{\rm even}}.
\]
显然，当改变局部坐标时，这个表达式不变。这给出超流形上一个整体定义的实向量场 $Q$，它是局部超对称生成元 $Q^i_\alpha$ 的扭变形式的 $(0,0)$ 分量的整体化。正因为存在这样一个整体的、处处非零的、满足 $Q^2=0$ 的超对称生成元，我们才能在 $X$ 上做整体的拓扑量子场论。

我们还在 $X$ 上有一个整体定义的一形式 $K$，取值于超流形上的向量场。它由 $\bar Q_{j,\dot\alpha}$ 的扭变形式导出。在同一局部坐标下，它的公式是：
\[
K=\bigl(\partial_{i,{\rm odd}}+\theta^{\mathbf{R}}\partial_{i,{\rm even}}\bigr)\,dx^i.
\]
由公式可以直接看出，$K$ 在坐标变换下不变，从而是 $X$ 上一个定义良好的整体对象。

% [Tutor Note: 「显然，当改变局部坐标时这个表达式不变」是全书最需要停下想一想的地方之一。$Q$ 的不变性并不是逐项成立：坐标变换 $x\to x'$ 会同时混合 $TX$ 与 $\Lambda^0TX$ 两个奇方向，也会让 $\theta^i$ 按 $\partial x'^i/\partial x^j$ 变换；两项之所以能拼成一个不变量，恰恰是因为 $\theta^i\partial_{i,{\rm even}}$ 与 $\partial_{\mathbf{R}^{\rm odd}}$ 的变换互相补偿（这一结构在超空间里称为「分裂超流形上的向量场由 $X$ 上的向量场与一个函数决定」）。这一点对 $K$ 同样成立，而 $K$ 的系数里出现的是 $\theta^{\mathbf R}$（不是 $\theta^i$），这正是 $K$ 成为一形式的关键。此外，原文「处处非零（everywhere nonzero）」四字不可省：若 $Q$ 在某处为零，则该点附近不能再作 $\{Q,\cdot\}$ 的上同调论证，整套拓扑不变性证明即失效。]

由超 Poincaré 代数结构，或由上述显式公式，容易推出 $Q^2=0=\{K,K\}$。由上述公式直接计算，或由一般的超对称代数结构，还给出
\begin{equation}
\tag{19.1}
\label{eq:19-1}
\{Q,K\}=d,
\end{equation}
其中 $d$ 是 $X$ 上的外微分。

% ===== 原文 p.266 =====
事实上，由 Noether 定理，这些方程有局部（即未积分）的形式。在 Minkowski 空间中，平移生成元 $P_\mu$ 由对能动张量（stress tensor）$T$ 积分得到，超对称荷同样由对超流（supercurrent）$\mathcal{S}$ 积分得到；$\mathcal{S}$ 是一个取值于旋量的三形式（或者按对偶性说，一个向量场）：
\[
\begin{aligned}
P_\mu&=\int_\Gamma d^3x\,T_{\mu0},\\
\bar Q_{\dot\alpha}&=\int_\Gamma d^3x\,\mathcal{S}_{\dot\alpha0}.
\end{aligned}
\]
扭变之后，超流分解为若干部分，其中一部分可以看作取值于一形式的三形式；我们就简单地把它记为 $\mathcal{S}$。上面第二个方程在扭变理论中变为
\[
K=\int_\Gamma\mathcal{S}.
\]
方程 19.1 的未积分形式是
\[
\{Q,\mathcal{S}\}=T.
\]
这个方程非常重要，因为它是我们即将发展的形式体系具有拓扑不变性的一个自然证明的基础。

% [Tutor Note: 原文「it a three-form (or by duality a vector field) with values in spinors」脱一「is」字，应读作「it \emph{is} a three-form」，译文按正确读法处理。另有一处作者未写出的次数核对：在 Minkowski 空间中 $\mathcal{S}$ 是\emph{取值于旋量丛}的三形式，故 $\Gamma$（三维）上的积分把三个方向积掉、只留下一个自由指标，这正是 $P_\mu=\int d^3x\,T_{\mu0}$ 与 $\bar Q_{\dot\alpha}=\int d^3x\,\mathcal{S}_{\dot\alpha0}$ 都只积三个方向的原因。扭变之后其中一块成为「取值于\emph{一形式}的三形式」，总次数 $3+1=4$，恰与 $X$ 的维数相合，这才使它能与 $T$（扭变后同样是取值于一形式的三形式）逐点配对；也正是因为 $\Gamma$ 是三维的，对 $\mathcal{S}$ 积掉之后剩下的一形式才是 $K$。若忽略「取值于一形式」这一层，$\{Q,\mathcal{S}\}=T$ 两边就相差一次，整套扭变记号随之失效。]

像所有这类关系一样，公式 $\{Q,\mathcal{S}\}=T$ 可以读作关于 Hamilton 向量场、或关于它们所生成的变换的陈述。然而——对于我们应用于整体四维流形 $X$ 的目的而言——这个公式最有力量的读法如下：$\mathcal{S}$ 与 $T$ 是 $X$ 上的局部算符；$Q$ 是由 $X$ 构造出的分裂超流形的奇对称 $Q=\partial_{\mathbf{R}^{\rm odd}}+\theta^i\partial_{i,{\rm even}}$ 在 $X$ 上算符代数中诱导出的整体定义的自同构。在这样的符号解释下，命题 $\{Q,\mathcal{S}\}=T$ 才有意义。

\subsubsection{$X$ 上的拓扑量子场论}
% 原文 §19.2.2：Topological Quantum Field Theory on $X$

我们将考虑 $Q$-闭（$Q$-closed）算符之乘积的期望值。由于任何 $Q$-恰当（$Q$-exact）算符的期望值都是平凡的，由此可知：任意多个 $Q$-闭算符之乘积的期望值为零，只要其中至少有一个是 $Q$-恰当的。

$\{Q,T\}=0$ 这一事实意味着，当我们变动度规时，关联函数的变化等于插入一个位于 $Q$ 的像中的算符。由此可知，我们将要考虑的关联函数不依赖于度规。我们心目中的算符，是在这个更一般的规范理论框架中 Donaldson 类的对应物。$Q$-闭算符的关联函数与度规无关，是 Donaldson 不变量在度规改变下不变这一性质的推广。

% [Tutor Note: 先声明记号，否则下引一步不成立：本节所有 $\{,\}$ 都是\emph{分次}括号——两个奇对象取反对易子，一奇一偶取对易子。在此约定下，作者由 $\{Q,\mathcal{S}\}=T$ 直接跳到「$\{Q,T\}=0$」，省掉的是分次 Jacobi 恒等式：$\{Q,T\}=\{Q,\{Q,\mathcal{S}\}\}=\frac12[[Q,Q],\mathcal{S}]=[Q^2,\mathcal{S}]=0$，因为 $[Q,Q]=2Q^2=0$。这一步是整讲「拓扑不变性」证明的枢纽，绝非平凡恒等式：它要求 $Q$ 是\emph{整体}定义且满足 $Q^2=0$ 的（见 p.265），而非仅仅局部存在。此外「度量改变给出 $Q$ 的像中的算符」一句的真实含义是：能动张量与度规变分相伴，故 $\delta_g\langle\cdots\rangle=\langle T\,\cdots\rangle=\langle\{Q,\mathcal{S}\}\cdots\rangle$，再用 $Q$-闭性把 $\{Q,\mathcal{S}\}$ 挪到一边即得零——这与 p.266 开头「$Q$-恰当算符期望值平凡」是同一件事的两种说法。]

让我们回忆从取值于函数的局部算符出发、生成取值于更高次微分形式的算符的一般方案。这里我们把第 11 讲与第 14 讲的讨论与结果从二维推广到四维。

从规范群 $G$ 的李代数 $\mathfrak{g}$ 上一个次数为 $d$ 的多项式函数 $P$ 出发，它在伴随表示下不变。经典地，把它作用于主 $G$-丛上的任意联络 $A$ 时，产生一个次数为 $2d$ 的 Chern 形式（Chern form）$P(F_A)$。把它作用于

% ===== 原文 p.267 =====
流形 $X$ 与构型的规范等价类空间 $\mathcal{B}$ 之积上的万有丛（universal bundle）的万有联络上，便得到 $X\times\mathcal{B}$ 上一个次数为 $2d$ 的闭形式。回忆前两讲：$\mathcal{N}=2$ 超对称规范多重态的一部分，是一个 $\mathcal{N}=1$ 手征多重态，其中含有一个取值于伴随丛的标量场 $\phi$。在量子力学层面，把多项式函数 $P$ 作用于 $\phi$，产生理论中一个 $Q$-闭的局部算符 $\mathcal{O}_P^{(0)}$。反复作用 $K$ 给出算符 $\mathcal{O}_P^{(n)}=K^n\mathcal{O}_P^{(0)}$。这些是取值于 $X$ 上 $n$ 形式的局部算符。对一个 $n$-圈 $\Sigma^n\subset X$，算符
\[
\int_{\Sigma^n}\mathcal{O}_P^{(n)}
\]
相当于把 $X\times\mathcal{B}$ 上万有联络的 Chern 形式 $P(F_A)$ 与圈 $\Sigma^n$ 作斜积（slant product），以在 $\mathcal{B}$ 上产生一个次数为 $2d-n$ 的闭形式 $\mu_P(\Sigma)$。我们将会看到，期望值
\[
\Bigl\langle\prod_{i=1}^s\int_{\Sigma_i^{n_i}}\mathcal{O}_P^{(n_i)}\Bigr\rangle
\]
等于
\[
\sum_E\int_{\mathcal{M}_E}\mu_P(\Sigma_1^{n_1})\cup\cdots\cup\mu_P(\Sigma_s^{n_s}),
\]
其中 $E$ 取遍 $X$ 上主 $G$-丛的拓扑类型，而 $\mathcal{M}_E\subset\mathcal{B}_E$ 是经典解的模空间。当然，只有当 $\mathcal{M}_E$ 的虚拟维数（virtual dimension）等于 $\sum_{i=1}^{s}(2d-n_i)$ 时，这个积分才有意义。

% [Tutor Note: 此处「次数」的记账必须自己动手核一遍。$P$ 的次数是 $d$，故 $P(F_A)$ 是 $2d$ 次闭形式；与 $n$-圈作斜积消去 $n$ 次，于是 $\mu_P(\Sigma)$ 在 $\mathcal{B}$ 上是 $2d-n$ 次——这正是文中 $2d-n$ 的来源。另一方面，$\mathcal{O}_P^{(n)}=K^n\mathcal{O}_P^{(0)}$ 中每作用一次 $K$ 就把形式次数提高一次（$K$ 是一形式），故 $\mathcal{O}_P^{(n)}$ 是 $n$ 形式，在 $n$ 维圈上可积；在四维流形上 $n>4$ 时它自动为零，故实际只有 $n=0,1,2,3,4$ 五种非平凡情形（$n=4$ 时 $\int_X\mathcal{O}_P^{(4)}$ 就是 $\int_X P(F_A)$ 型的整体不变量）。把所有插入合起来：$\bigcup_i\mu_P(\Sigma_i^{n_i})$ 的次数是 $\sum_i(2d-n_i)$，只有它恰等于 $\mathcal{M}_E$ 的虚拟维数时，$\int_{\mathcal{M}_E}$ 才可能不为零——这便是「虚拟维数等于 $\sum(2d-n_i)$ 时该积分才有意义」的确切含义。]

由 $\{Q,K\}=d$ 以及 $\{Q,\mathcal{O}^{(0)}\}=0$，我们有
\[
\{Q,\mathcal{O}^{(n)}\}=d\mathcal{O}^{(n-1)}.
\]
这意味着 $\mathcal{O}^{(n)}$ 在 $n$-圈上的积分是 $Q$-不变的。

\subsubsection{$\mathcal{N}=2$ 矢量多重态的扭变与整体化}
% 原文 §19.2.3：Twisting and globalizing the $N=2$ vector multiplet

在能把超对称理论安置到一般的 Riemann 流形上之前，我们还有工作要做。到目前为止，我们研究了扭变对超欧氏群（super-Euclidean group）的作用，并看到如何在定向 Riemann 流形上实现它。我们仍须考察它对 Lagrange 量中诸场的影响。本节研究矢量多重态。

让我们用 $\mathcal{N}=1$ 的语言回忆 $\mathcal{N}=2$ 矢量多重态中（未扭变的）场的基本构型。首先，$\mathcal{N}=2$ 矢量多重态分解为一个 $\mathcal{N}=1$ 矢量多重态
\[
\mathcal{A}=(A,\lambda,D)
\]
与一个 $\mathcal{N}=1$ 手征多重态
\[
\Phi=(\phi,\zeta,F).
\]
在第一个（矢量）多重态中，$A$ 是规范场，$\lambda$ 是取值于伴随丛的旋量，$D$ 是取值于伴随丛的实辅助场（auxiliary field）。在手征多重态中，$\phi$ 是伴随丛复化的一处截面，$\zeta$ 是取值于伴随丛的旋量，$F$ 是取值于伴随丛的复值辅助场。在 $SU(2)_R$ 对称性下，对 $(\lambda,\zeta)$ 按一个二维表示变换，该表示记为 $\psi^i_\alpha$；而 $\vec D=(D,\operatorname{Re}F,\operatorname{Im}F)$ 按 $SU(2)_R$ 的伴随表示变换。

% ===== 原文 p.268 =====
下表归纳了各种场的 $U(1)_R$ 荷。

\begin{center}
$\mathcal{N}=2$ 矢量多重态（The $\mathcal{N}=2$ Vector Multiplet）\\[4pt]
\begin{tabular}{|c||c|c|c|c|c|}
\hline
$U(1)_R$ 荷 & $-2$ & $-1$ & $0$ & $1$ & $2$\\
\hline
场 & $\bar\phi$ & $\bar\psi_{\dot\alpha j}$ & $A$, $\vec D$ & $\psi^i_\alpha$ & $\phi$\\
\hline
\end{tabular}
\end{center}

现在按上述方式用 $SU(2)_R$ 作扭变。$SU(2)_R$ 除在 $\vec D$ 与费米子上之外作用平凡；因此这些是扭变时唯一改变性质的场。扭变之后，$\psi^i_\alpha$ 按 $(\frac12,\frac12)$ 表示变换。这意味着，扭变之后 $\psi^i_\alpha$ 成为取值于伴随表示的一形式；我们把它记为 $\psi_{(1)}$。类似地，扭变之后 $\bar\psi_{\alpha j}$ 按 $(\frac12\otimes\frac12,0)$ 变换，从而分解为一个函数 $\psi_{(0)}$ 与一个自对偶二形式 $\psi_{(2)}$ 的直和。辅助场 $\vec D$ 也按 $(1,0)$ 变换，从而成为 $X$ 上的一个自对偶二形式。

我们看到，扭变之后所有这些场都在 $X$ 上整体存在，因为它们已成为微分形式。这意味着，就纯规范理论而言，我们可以在任意闭定向四维流形上写下保持超对称荷 $Q$ 的 Lagrange 量。

% [Tutor Note: 正文此处把该费米子写作 $\bar\psi_{\alpha j}$，而上表及前文一律作 $\bar\psi_{\dot\alpha j}$（$\alpha$ 上带点）。这是原文的记号不一致，并非两种场：$\bar\psi$ 的旋量指标必须带点（属 $SU(2)_-$），否则它不可能按 $(\frac12\otimes\frac12,0)$ 变换——带点的 $\dot\alpha$ 提供 $SU(2)_-$ 的 $\frac12$，$j$ 提供 $SU(2)_R$ 的 $\frac12$，对角化后才是 $0\oplus1$，从而分解出函数 $\psi_{(0)}$ 与自对偶二形式 $\psi_{(2)}$。另可顺带核对荷的账：由上式 $\phi$ 的 $U(1)_R$ 荷为 $2$，故 $P(\phi)=\operatorname{Tr}\phi^2$ 的荷为 $2\times2=4$，恰好等于它在下文所对应的 Donaldson 类 $\mu(\text{pt})\in H^4(X)$ 的次数——作者正是靠这个对应确认 $\operatorname{Tr}\phi^2$ 就是 Donaldson 理论中那个四次类的对应物。这也是核对荷表最简便的办法：若把 $\phi$ 的荷误记成 $1$，$\operatorname{Tr}\phi^2$ 就只有荷 $2$，与 $H^4$ 对不上。]

让我们考察 $Q$ 对这些场的作用。场 $(A,\psi_{(1)},\phi)$ 连同作为微分的 $Q$，构成规范群作用在联络空间上的等变上同调（equivariant cohomology）的一个模型。为了使记号更像一个微分，我们以 $\delta$ 记 $Q$ 的作用。于是有
\[
\begin{aligned}
\delta A&=\psi_{(1)},\\
\delta\psi_{(1)}&=-d_A\phi,\\
\delta\phi&=0.
\end{aligned}
\]
只用这些场时，$Q$-闭算符有模 $Q$-恰当算符的代表元。我们从一个由 $\phi$ 的规范不变多项式给出的算符 $\mathcal{O}_P^{(0)}=P(\phi)$ 出发。例如在 $SU(2)$ 情形，取 $P(\phi)=\operatorname{Tr}\phi^2$。该多项式具有 $U(1)_R$ 荷四，对应于第四上同调中的 Donaldson 类 $\mu(\text{pt})$。由这些局部算符，我们按公式 $\mathcal{O}_P^{(j)}=K^j\mathcal{O}_P^{(0)}$ 造出取值于形式的算符。

$\delta$ 对其余场的作用是：
\[
\begin{aligned}
\delta\bar\phi&=\psi_{(0)},\\
\delta\psi_{(0)}&=[\phi,\bar\phi],\\
\delta\psi_{(2)}&=F_A^+-\vec D,\\
\delta(F_A^+-\vec D)&=[\phi,\psi_{(2)}].
\end{aligned}
\]
这里我们依照的记号是：曲率 $F_A$ 分解为自对偶与反自对偶两部分 $F_A^++F_A^-$。其解释如下。若只纳入 $\bar\phi$ 与 $\psi_{(0)}$，就得到等变上同调的一个模型，所涉及的是规范

% 原文 p.268 至此结束。末句断在「…the equivariant cohomology of the gauge」处，全句为
% 「If one includes just $\bar\phi$ and $\psi_{(0)}$, one gets a model of the equivariant
% cohomology of the gauge group that was described in Lecture 14.4.」，其余部分（「group that was
% described in Lecture 14.4」）在 p.269，由 lec19_p2 以「……第 14.4 讲所描述的那个群」续起。
% 拼合时本页句末的「规范」与下页的「群」直接相接即可。
